% !PRINTED 599
\weberpage{599}
%<<CONT
ob eine oder einige der Wurzeln auf diese Weise dargestellt
werden können, während andere eine solche Darstellung nicht
gestatten. Diese Frage ist nur berechtigt bei irreduciblen Gleichungen, da bei reduciblen Gleichungen alle denkbaren Combinationen vorkommen können, und hier gilt nun der Satz:

\begin{quote}
\noindent IV.\ \gesperrt{Wenn eine Wurzel einer irreduciblen Gleichung
durch Lösung cyklischer Gleichungen bestimmbar ist, so ist die Gleichung metacyklisch.}
\end{quote}

Wenn eine irreducible Gleichung $n^{\text{ten}}$ Grades $f(x) = 0$ auch
nur eine Wurzel hat, die durch successive Adjunction von Wurzeln
cyklischer Gleichungen rational wird, so muss sie nothwendig
durch diese Adjunction reducibel werden, da sich ja schliesslich
ein linearer Factor absondern muss. Es sei also $P$ die Gruppe
unserer Gleichung, nachdem alles Nöthige so weit adjungirt ist,
dass zwar $f(x)$ noch nicht reducibel ist, aber durch die nächste
Adjunction der Wurzel $\varepsilon$ einer cyklischen Gleichung von Primzahlgrad $m$ in Factoren zerfällt. Es muss dann $P$ nach \S.\,175
einen Normaltheiler $Q$ vom Index $m$ haben, auf den sich die
Gruppe der Gleichung nach Adjunction von $\varepsilon$ reducirt, und die
Permutationsgruppe $Q$ muss intransitiv sein. Wenn $Q$ die Einheitsgruppe ist, so ist $f(x) = 0$ durch Adjunction von $\varepsilon$ vollständig gelöst. Ist aber $Q$ noch von der Einheitsgruppe verschieden, so treten die Sätze des \S.\,158 in Kraft. Danach zerfällt
$f(x)$ nach Adjunction von $\varepsilon$, da ja hier der Index von $Q$ eine
Primzahl ist, in $m$ Factoren, $\varphi(x, \varepsilon), \varphi_1(x, \varepsilon), \ldots \varphi_{m-1}(x, \varepsilon)$, vom
gleichen Grade $\mu$ und es ist $\mu m = n$.

Ist z. B. $n$ eine Primzahl, so muss $\mu = 1$ sein; die Functionen $\varphi$ sind linear, und die Gleichung $f(x) = 0$ ist vollständig gelöst; also ist der Satz IV für Gleichungen von Primzahlgrad richtig. Im allgemeinen Falle muss einer der Factoren
$\mu^{\text{ten}}$ Grades $\varphi(x, \varepsilon), \varphi_1(x, \varepsilon), \ldots \varphi_{m-1}(x, \varepsilon)$ etwa $\varphi(x, \varepsilon)$ eine
Wurzel haben, die durch Adjunction der Wurzeln cyklischer
Gleichungen rational wird, und wenn wir also annehmen, dass
unser Satz für Gleichungen $\mu^{\text{ten}}$ Grades schon bewiesen sei, so
folgt, dass $\varphi(x, \varepsilon) = 0$ selbst und also auch ihre Gruppe metacyklisch ist.

Da aber nach \S.\,158 die verschiedenen Gleichungen $\varphi = 0$,
$\varphi_1 = 0, \ldots \varphi_{m-1} = 0$ dieselbe Gruppe haben, so sind sie alle
und mithin auch $f(x) = 0$ metacyklisch.
